%--
\chapter{AI时代的算法课}
\begin{flushright}
    —— \textit{人工智能时代，算法课还需要教什么，应该怎么教？}
\end{flushright}

{\itshape
    按：人工智能能力日进，“算法设计与分析”课程，无论学生还是教师，都免不了被追问或自问一个问题：课堂上讲授的这些基础知识，还有没有用？还要不要学？还要不要按传统课堂讲授的方式来教或者学？我从授课教师的角度，把自己的思考整理成文。全文的基本论断只有一句：简单的知识传授，在人工智能时代已经不再有意义；算法课着力讲授的应该是知识点之间、之上、之外的隐知识。
}

%----------------------------------------
%\section{基本论断——知识点之间、之上、之外}

\medskip

即使在人工智能出现之前，经典算法知识的学习资料就早已极大丰富——经典教材、经典授课视频，以及包括习题与答案在内的各种配套资料，应有尽有。人工智能时代，这类知识的获取与简单加工变得更加容易。因此，“把知识点本身讲给学生”这件事，价值已经不大。
算法课着力讲授的应该是知识点之间、之上、之外的隐知识。这是全文立论的基础。
围绕这一论断，我们分三部分展开论述。

%----------------------------------------
\section{知识点“之间”与“之上”}

\nid\textit{关键词：索引、按需建索引；抽象、共性、抽象层级}

\medskip

第一层面的“之间”与“之上”大致对应\textit{广度}与\textit{深度}这两种视角。

%--------------------
\subsection{之间：知识的梳理与索引}

从广度的视角看，算法课的知识点组织，本身就像是一个\textit{索引}。
在简单的知识传递之上，我们的授课过程，就像是数据库建立索引的过程。
学生未来必将面对远远超出其所能记住、所能深知的知识体量。
面对这样大的知识体量，重要的是两种能力：
%--
\begin{enumerate}[label=\arabic*.]
    \item 对已经学到、或必须要学的“内核级”知识，脑中有一个很好的\textit{索引}；
    \item 对未来要学的更多知识，能够\textit{按需建立索引}。
\end{enumerate}

\nid 这种索引式的知识组织能力，正是算法课的讲授希望教授给学生的。

%--------------------
\subsection{之上：抽象的能力}

从深度的视角看，在具体的知识点之上，可以从问题、策略、数据结构、数学结构等不同的角度进行抽象；而且抽象的层级可以逐层多次提升，相应所捕捉到的共性也就越来越深刻。

这种抽象能力的价值在于：未来面对所谓的新场景、新问题，学生能够经过适当的\textit{抽象层级的提升}，捕捉到必要的共性，从而有效地基于已经学到的知识，解决新场景中的新问题。

\medskip
\nid\textit{小结：第一层面，广度上是“建索引”的组织能力，深度上是“提层级”的抽象能力。二者分别在知识点之间、之上，却又都以知识点为依托。}

%----------------------------------------
\section{知识点“之外”}

\noindent\textit{关键词：创新、灵光乍现、熟练掌握、跳板；统觉的综合统一、图式}

\medskip
在第一层面之上，还有第二个层面：\textit{创新}。

%--------------------
\subsection{表面性的阐述}

对“创新”这个层面，有一种看法认为：只要你的抽象层级足够高，那就没有什么事情是“新”的，总都有某种共性。照此说来，“创新”这个层面是不是根本就不应该存在？

回顾实践的过程，我认为上述提法虽有一定道理，但是在实践中可行性不足。当抽象层级提得足够高，“面对真实问题—提升抽象层级—共性涌现—解决问题”这个流程，很多时候并不能顺畅地完成。真实场景中的真实问题，往往还需要\textit{灵光乍现}的过程。

如何理解这个灵光乍现的过程，是一个深刻而困难的问题。从外在表象上说：经过第一层面充分多的训练，以既有知识的熟练掌握为\textit{跳板}，我们可能获得更多的创新能力——面对没有简单共性的新场景、新问题，能够拿出好的解决办法。但是，在这一表象的背后，到底发生了什么。我（和我的合作者）在2009年的一篇学术论文中引用过康德的先验哲学理论，放在这里或许可以对难以言说的“灵光乍现”，提供一种有益的解释。

%--------------------
\subsection{康德先验哲学视角下的“灵光乍现”}

上面“以熟练掌握为跳板”的朴素表述，其实触及了一个古老的哲学问题：人的认识能力如何在教育与训练中形成？用康德的先验哲学来解读，这段话可以获得一种更精确、更深刻的阐述。

%----------
\subsubsection{（1）“灵光乍现”的本质：先验综合对杂多的重构}

康德对洛克经验论的核心批判是：\textit{经验只能给予材料，不能给予形式}。放到算法学习的场景中，一个学生面对一道全新的习题或开放性的算法问题，他给出的解法里必然包含两个成分：
%--
\begin{itemize}
    \item \textit{经验性成分}：课堂上学过的排序、最短路径、分治、贪心、动态规划等具体知识，做过的习题——这对应康德所说的感性\textit{杂多}（Mannigfaltige）；
    \item \textit{先验性成分}：那种使他能把这个新问题“把握为一个问题”、能在头脑中组织、变换、重组条件与目标的统一性活动——这在康德那里对应着时间形式、范畴，以及最根本的，\textit{统觉的综合统一}。
\end{itemize}

“灵光乍现”之所以显得神秘，是因为我们习惯把它理解为经验的奇迹——仿佛新解法是从旧习题里“长”出来的。但康德会说：真正闪现的从来不是旧知识本身，而是\textit{综合活动本身的展现}。新解法的“新”，在于主体的综合结构对全新杂多的重新立法。康德的名言“\textit{思维无内容则空，直观无概念则盲}”在这里精确成立：不经过第一层面的积累，形式空转；而只有积累、没有内在的综合能力，题目就永远只是背过的题型，而不是一个需要被理解的问题。


%----------
\subsubsection{（2）“以既有知识的熟练掌握为跳板”：三重综合的训练学意义}

康德在“三重综合”中的分析，可以用来精确重述“跳板”这个比喻。跳板不是知识量的堆积，而是\textit{综合能力的锤炼}：
%--
\begin{itemize}
    \item \textit{直观中领会的综合}：学生做题时，把题目条件逐条把握为一个整体——“给定”“求”“约束”在意识中被统一起来。这本身就是需要训练的能力，初学者常常读不懂题，正是因为杂多还没有被领会为统一的问题结构。
    \item \textit{想象中再生的综合}：调用递归式、分治法、不变式等知识点——这些表象能在意识中被重新唤起并联结。康德在“经验的类推”中指出，表象间的习惯性联想本身就是想象力综合的产物。所谓“熟练掌握”，按康德的解释，就是\textit{让算法设计的表象之间建立起稳定的再生联结}：看到“最优子结构”，动态规划的表象自动浮现；看到“重叠子问题”，暂存子问题结果的表象自动跟进。这种联想网络越致密，跳板就越结实。
    \item \textit{概念中认定的综合}：把题目中的具体情境认作“这是图遍历问题”“这是可以用不变式证明算法正确性的问题”。这一步最关键——它正是康德所说的\textit{用概念赋予杂多以统一性}。
\end{itemize}

因此，“充分多的训练”在康德式解读下，不是刷题的机械重复，而是\textit{想象力图式（Schemata）的锻造}。康德认为概念与直观之间必须有图式这个中介——抽象规则本身不能直接应用到具体问题上，必须借助图式。算法教育的目标，本质上就是让学生\textit{内化}大量图式：把“不变式”这个抽象概念，通过归纳法证明的图式，锚定在可操作的直觉上。训练的积累之所以能成为跳板，是因为图式库足够丰富之后，新场景便不再完全陌生，而是可以被既有图式重新编码、再加以改造的对象。


%----------
\subsubsection{（3）“没有简单共性的新场景”：范畴对共性的超越}

“没有简单共性”是这段话的关键，也正是康德哲学最能照亮之处。朴素的学习观隐含一个假设：解决新问题靠的是新旧问题之间的相似性——创新是“迁移”。但如果新场景真的没有简单共性，迁移何以可能？

康德的回答是：\textit{我们能把全新经验把握为“一个经验”，靠的不是表层的共性，而是深层的范畴结构}。因果性（这个问题的结果由哪些条件制约）、量（解的规模、复杂度的高低）、模态（这个问题是否必然可解、在什么条件下可能可解）——这些不是从任何具体题目中归纳出来的，而是任何问题被思考为问题时就已经动用的形式。因此，面对全新问题仍能拿出好办法的学生，其行为印证了康德的核心主张：\textit{对象必须符合主体的认识形式}。他不是在新问题里“认出”了旧题的影子，而是用自己内化的综合结构去\textit{重构}了问题本身。所谓创新能力，不是发现新问题里藏着旧知识，而是\textit{先验结构对新杂多的立法}——这正是“人为自然立法”在学习者身上的体现：一个训练有素的学生所面对的“问题”，早已不是题目原文的杂多，而是经过他的认识形式整理过的、结构化了的问题。

%----------
\subsubsection{（4）康德式的边界提醒：对“灵光乍现”本身保持批判的谦逊}

康德严格区分\textit{现象}与\textit{物自体}，这对算法教育有一个冷静而深刻的提醒：
%--
\begin{itemize}
    \item 我们从外部能观察到的，是现象层面：学生刷了多少题、考试分数的跃迁、某天突然解出了难题；
    \item 而“为什么经过这些训练综合统一能力就被激发出来了”——这触及的是认知过程自身的“物自体”。
\end{itemize}

康德的教训是：\textit{对于物自体，理性合法的言说极为有限}。因此，一个优秀的算法教师面对“学生什么时候开窍”这个问题，应当保持一种康德式的谦逊——正如上文所说“还没有特别深刻的认识，只能从外在表象上说”。这不是认识论上的失败，而是认识成熟的标志。康德说批判的目的不是扩展知识，而是\textit{划定知识的界限}。承认教育机制中的不可完全解释之处，恰恰是把教学当作一门严肃学问而非手艺的前提。反过来，任何声称“我已经完全掌握了创新能力如何产生”的教学理论，在康德看来都需要警惕——当理论越过经验所能担保的范围，就会陷入先验幻相。


%----------
\subsubsection{（5）一点延伸：先验与经验的张力}

最后值得注意的是康德框架与本文论点之间真正深刻的张力。康德本人坚持先验形式是先天的、不可由经验生成的。但算法训练的经验似乎暗示了一种反向关系：\textit{经验的积累塑造并强化了形式}——综合规则、图式网络、问题图式的重构方式，确实是在大量做题与纠错中被锻造出来的。这正是康德框架在此受到的挑战：如果综合的统一能力可以经由训练被发展、被改造，那么“先验”就失去了绝对的先天性，成为一种\textit{历史地获得的先验}——更像皮亚杰或发生认识论所说的、内化的操作结构。康德未必接受这一修正，但正因如此，用他的框架对照本文，我们既能看清“灵光乍现”是一个古老的哲学问题，也能看清算法教育把它推进到了哪里：教育的问题，归根结底是\textit{如何为学生的认识能力立法}的问题——不是灌输内容，而是锻造形式。

\medskip

\nid\textit{小结：创新也许无法直接传授，但可以经由充分的第一层面训练而被“诱出”；熟练掌握，是灵感的跳板。用康德的语言重述即是：充分的习题与知识积累，使学生本已具备的综合统一能力获得了充分的表现；创新能力不是知识量的产物，而是认识形式显现自身的过程；我们只能从学生的解题行为中推断这种能力的存在，却无法直观其内在机理——对这种不可知的自觉承认，恰恰是批判哲学所能达到的最高诚实，也是一切认真面对教育之神秘性的教师应有的姿态。}

%----------------------------------------
\section{隐知识视角下的算法课堂教学}


\noindent\textit{关键词：隐知识（tacit knowledge）、默会认识、元知识、高带宽的非语言通道、物理共处}

前两节论证了算法课应该\textit{教什么}：知识点之间的索引、知识点之上的抽象、知识点之外的创新。本节换一个视角，回答按语中的另一半问题——\textit{还要不要按传统课堂黑板教学的方式来讲授}。所用的理论工具，是波兰尼（Michael Polanyi）的隐知识学说。

%--------------------
\subsection{隐知识的核心命题：认识总有一个不可完全对象化的维度}

波兰尼所说的，不只是“有些知识暂时没写出来”。他的命题要深刻得多：\textit{任何认识都包含一个不能被完全对象化的默会维度}——默会认识是一种整合线索、形成意义的认识\textit{活动}，而不只是头脑中储存的隐性\textit{内容}。

这个视角恰好与上一节的康德式解读互补。面对“灵光乍现”，康德的框架提醒我们保持谦逊：其内在机理如物自体，不可直观、不可尽言。波兰尼则给出更积极的半句：这个维度虽然不可完全\textit{对象化}——无法被完整写成条文、编进教材——却并非不可习得；它是可以在认识活动中被亲身\textit{践行和锻炼}的。不可言说，不等于不可获得。

%--------------------
\subsection{课堂的不可替代性：隐知识依赖高带宽的非语言通道}

隐知识如何习得？这里需要引入一个相关的概念：元知识——关于“如何获得、提炼、调用知识”的知识。元知识的传播，依赖\textit{高带宽的非语言通道}。学习者获得这三种能力（获得、提炼、调用）中的相当一部分，是通过\textit{观察一个讲授者如何思考所遇到的困难、如何放弃一个既定策略、如何修改或者替换方案}而习得的。这种信息无法被文本载体编码，只能即时交互、即时反馈的环境中“泄露”出来，而（物理的）课堂，正是学生可以低成本获得的这样的一种环境。这就是为什么远程教育在传授显式的、具体的知识内容时高效，传授隐式的知识时近乎失效。

由此回看算法课，AI 时代的分工就清晰了。教材、授课视频、大模型，都在\textit{显知识}的一侧——它们传递知识点高效而廉价，这正是前文“简单知识传授不再有意义”的原因。而本文所论的索引式组织、抽象层级提升乃至创新，很大程度上属于\textit{元知识与隐知识}的一侧：它们无法被完全编码，只能在师生共处的课堂中，通过观察教师真实的思维过程来习得。

这就回答了按语中的问题：\textit{算法课堂仍然要教，但教的重心必须移向隐知识}——黑板前的推演、（精心设计的）失败与重试，不是教学的冗余，恰恰是教材和大模型无法替代的部分。

%--------------------
\subsection{隐知识的可教性：演示思维、促成反思、以实践固之}

隐知识并非全然神秘，它的可教性有一条清晰的谱系。高度情境化、尚未充分外化的元知识，其获得往往依赖观察、模仿和共同实践；其中\textit{一部分}可以通过反思和案例比较被提炼为显性规则——但规则即便写下，能否在适当情境中被\textit{调用}，仍取决于实践经验、反馈和持续的自我监控。

对算法课堂教学而言，这意味着三重任务：
%--
\begin{enumerate}[label=\arabic*.]
    \item \textit{演示思维，而非只展示答案。} 教师应把卡壳、试错、改进的过程暴露在课堂上——学生从中“泄露”式地习得的，正是教科书里不存在的东西；
    \item \textit{促成反思，把默会经验部分外化。} 通过案例比较、回顾复盘——这正是我们的tutorial课程所强调的——帮助学生把散落的默会经验提炼为显性规则，编入他们在第一层面建立的“索引”；
    \item \textit{以实践与反馈固之。} 规则的可调用性只能靠实践经验、反馈与持续的自我监控来养成——习题课、讨论与反复试错，因此不是讲授的附属，而是隐知识传递的主场。
\end{enumerate}

\nid\textit{小结：康德说“灵光乍现”的机理不可尽言，波兰尼补上另一半：不可完全对象化的东西，仍可在共处与实践中习得。AI 接管了显知识的传递，恰恰把课堂还给了它最本己的任务——高带宽、低延迟、物理共处之中的隐知识传递。算法课要教的，不只是索引、抽象与创新这些“内容”，更是让学生在与教师真实的思维相遇中，获得那些无法被文本编码的能力。}

%----------------------------------------
\section{总结——从教什么到怎么教}

把全文的线索串起来，结论如下：
%--
\begin{enumerate}[label=\arabic*.]
    \item 人工智能时代，知识点本身的传授已不再稀缺；算法课的作用，在知识点\textit{之间、之上、之外}；
    \item 知识点\textit{之间}与\textit{之上}：广度上建立\textit{索引式}的知识组织，深度上训练\textit{抽象共性、提升层级}的能力——面向“有共性可循”的新问题；
    \item 知识点\textit{之外}：\textit{创新}——面向“没有简单共性”的新问题。它未必可教，但可以经由充分的第一层面训练、以熟练掌握为跳板而获得；其机理可用康德先验哲学加以阐发，同时应对其保持批判的谦逊；
    \item 至于\textit{怎么教}：索引、抽象与创新都含有不可完全编码的\textit{隐知识}维度，其传递依赖高带宽的非语言通道与物理共处。AI 接管显知识的传递之后，课堂的重心必须移向隐知识——演示思维、促成反思、以实践与反馈固之。
\end{enumerate}


三个维度阐明了算法课\textit{教什么}，如表\ref{tab:three-dimensions-algorithm-knowledge}所示；而隐知识视角进一步阐明\textit{怎么教}——这三个维度中的每一项，都有一个无法被文本完全编码的默会维度，只能在师生共处的课堂中，经由观察、模仿、实践与反馈来习得。显知识的传递交给书籍与大模型，隐知识的传递还给课堂——这正是当前算法课讲授所遵循的逻辑。

%--
\begin{table}[htbp]
    \centering
    \caption{算法课着力讲授的内容}
    \label{tab:three-dimensions-algorithm-knowledge}
    \begin{tabularx}{\textwidth}{l l >{\raggedright\arraybackslash}X >{\raggedright\arraybackslash}X}
        \toprule
        \textbf{维度} & \textbf{视角} & \textbf{培养的能力} & \textbf{面向的问题} \\
        \midrule
        知识点\textit{之间} & 广度 & 索引式的知识组织 & 知识体量远超记忆与掌握 \\
        知识点\textit{之上} & 深度 & 抽象共性、提升抽象层级 & 有共性可循的新场景、新问题 \\
        知识点\textit{之外} & 跃迁 & 创新（灵光乍现） & 没有简单共性的新场景、新问题 \\
        \bottomrule
    \end{tabularx}
\end{table}

%----------------------------------------
\section{余论——未尽之意与待答之问}

上面的论述，更多是在阐述算法课教什么、怎么教，但并未阐明人工智能时代，所教的上述知识是否有用，如何运用。
这一内容留待下一篇文章的阐述，并且鉴于技术与时代的迅速进步，所需的可能是“西西弗斯”式的不断地思考与阐述。

%----------
\subsubsection{进一步阅读}

%--
\begin{enumerate}
    \item 康德. 纯粹理性批判[M]. 邓晓芒, 译. 北京: 人民出版社, 2004.
    \item 洛克. 人类理解论[M]. 关文运, 译. 北京: 商务印书馆, 2011.
    \item 皮亚杰. 发生认识论原理[M]. 王宪钿, 等译. 北京: 商务印书馆, 2017.
    \item Michael Polanyi, The Tacit Dimension, Doubleday \& Company, INC., Garden City, New York, 1966.
\end{enumerate}
